\reading{IM-8}{Hedge Fund Investment Strategies}{STRIKE FIRST}{3}

\srcnote{David P. Stowell and Paul Stowell, \textit{Investment Banks, Hedge
Funds, and Private Equity}, 4th Edition (Elsevier, 2023), Chapter 12.}

\begin{imtrapbox}[breakable=false,title={No GARP source exists for this reading}]
IM-8 is one of eight Investment Management readings with \emph{no GARP source
book} in the curriculum set. Where this volume fills a gap or resolves a
question, it does so from the available \term{Schweser condensation} --- one rung
below a GARP source. Schweser's framing may compress or reword GARP's, and no
Schweser wording is presented here as GARP wording. Every such passage is badged.
\end{imtrapbox}

\begin{imkeybox}[title={The four objectives in this reading}]
\begin{enumerate}[label=\textbf{\alph*.}]
\item Describe and compare different types of equity-based and macro hedge fund strategies and explain their return and risk profiles.
\item Compare different types of arbitrage strategies and describe the mechanics of convertible arbitrage strategies.
\item Describe the mechanics and risks of event-driven strategies such as merger arbitrage.
\item Explain the mechanics of a distressed securities strategy, and describe challenges, risks and considerations that may arise when implementing this strategy.
\end{enumerate}
\end{imkeybox}

% ------------------------------------------------------------------
\lo{IM-8 a}{Describe and compare different types of equity-based and macro hedge fund strategies and explain their return and risk profiles.}

\subsection{Equity long/short}

An equity long/short strategy uses fundamental analysis to identify undervalued
(long) and overvalued (short) stocks, aiming to profit from mispricing while
reducing overall market risk. Managers take positions based on industry,
geography and macro factors, with a defined investment horizon for price
convergence. For example, a fund may go long potential takeover targets while
shorting weaker firms in the same sector.

\term{Timing is crucial, especially for short positions}, as rising markets can
lead to margin calls or short squeezes. To control risk, funds diversify short
positions so that no single position dominates, and use factor models to reduce
exposures beyond market beta.

\begin{imkeybox}[title={Sources of return, and the costs against them}]
\textbf{Returns:} alpha from both long \emph{and} short positions through correct
security selection; the \term{interest rebate} earned on proceeds from short
selling, after deducting borrowing costs; and interest earned on cash held as a
liquidity buffer for margin requirements.

\smallskip
\textbf{Costs:} financing costs on margin, stock borrowing fees, and
\term{dividends paid on shorted shares}; plus transaction costs from frequent
trading and rebalancing.
\end{imkeybox}

\subsection{Dedicated short (short-biased)}

A dedicated short strategy focuses on identifying overpriced securities and
selling them short to profit from price declines. It involves high risk due to
margin calls, short squeezes and forced position closures.

Since most mutual funds focus on buying undervalued stocks and rarely short,
\term{overpriced securities are often less analysed}, creating opportunities for
short specialists. These managers look for signs like aggressive accounting,
legal risks or misleading disclosures. However, because markets tend to rise over
time, short positions face a \term{natural headwind}, making this strategy
difficult without strong expertise.

\subsection{Nonhedged equity}

A nonhedged equity strategy also relies on fundamental analysis but is
\term{long only}, with a focus on individual firms or industries. It resembles a
more traditional stock-picking approach: concentrated bets, with both upside and
downside analysed more deeply than a mutual fund would. Returns are enhanced
through \term{leverage} to magnify small mispricings, and through structures like
bull call spreads that profit from moderate price increases.

\subsection{Macro strategies}

\begin{itemize}
\item \term{Global macro.} Large, often leveraged positions based on expected
movements in broad markets --- equities, fixed income, interest rates, foreign
exchange and commodities --- using options, futures and swaps across asset
classes, with the flexibility to invest globally without restrictions. Managers
allocate across regions and asset classes based on economic outlook, interest
rate movements, capital flows, political events and policy changes. These
strategies typically focus on \term{highly liquid instruments}.
\item \term{Emerging markets.} Investing in developing economies through equities
and sovereign or corporate debt. High growth potential from industrialisation and
economic expansion, against higher political instability, currency volatility and
weaker regulatory frameworks.
\end{itemize}

\begin{imtrapbox}[breakable=false,title={Liquidity runs the opposite way in the two macro strategies}]
Global macro concentrates on the \emph{most} liquid instruments in the market;
emerging markets accepts some of the least liquid. Both are ``macro'' and both
are directional, so a question that swaps their liquidity profiles is using the
only structural difference between them.
\end{imtrapbox}

% ------------------------------------------------------------------
\lo{IM-8 b}{Compare different types of arbitrage strategies and describe the mechanics of convertible arbitrage strategies.}

Arbitrage strategies aim to exploit pricing inefficiencies. These occur when the
same asset trades at different prices in different markets, when two assets with
similar cash flows are priced differently, or when an asset with a known future
price is not priced at its risk-free discounted value. The core idea in every
case is: \term{buy the undervalued asset, sell the overvalued asset, profit when
prices converge.}

\subsection{Fixed-income arbitrage}

A long/short approach targeting pricing inefficiencies among fixed-income
securities. The core idea is that prices or yields of similar bonds should
converge over time: \term{low risk per trade, small but consistent returns}.

\begin{itemize}
\item Long Italian bonds (hedged) and short US Treasuries if the yield spread is
mispriced.
\item Short the expensive \term{on-the-run} bond and long the cheaper
\term{off-the-run} bond, as the on-the-run premium fades over time.
\end{itemize}

\subsection{Convertible arbitrage}

\gapbadge
\srcnote{The source notes carry the heading ``2) Convertible Arbitrage'' on an
otherwise blank page --- the mechanics the objective names explicitly are absent.
Written below from the available Schweser condensation, \term{one rung below a
GARP source}; no GARP text exists for this reading.}

\begin{imgapbox}[title={What the objective asks for}]
Not a definition of a convertible bond, but the \emph{mechanics of the trade}:
what is bought, what is shorted, why the hedge has to be rebalanced, and where
the money actually comes from.
\end{imgapbox}

A convertible bond contains an \term{embedded equity call option}. The bondholder
has the right to convert into a specified number of common shares and will
exercise when the value of the shares received exceeds the bond's cash redemption
value. Until conversion, the investor receives coupons like any bondholder.

\begin{imkeybox}[title={Three sensitivities, three value components}]
A convertible bond's price is sensitive to changes in \term{interest rates}, the
issuer's \term{credit risk}, and movements in the issuer's \term{share price and
price volatility}. Correspondingly, three components contribute to its value:
\term{investment value}, \term{conversion value} and \term{option value}.
\end{imkeybox}

\subsubsection*{The trade}

Convertible bond arbitrage involves \term{purchasing a convertible bond and
simultaneously selling short the underlying stock}. The arbitrageur seeks to
profit from both the bond's income and equity volatility. Although the hedge is
designed to neutralise equity risk, it must be \term{actively managed}: changes
in the stock price alter the delta of the convertible bond, requiring continuous
rebalancing, a process known as \term{delta hedging}.

\begin{itemize}
\item \textbf{Returns:} the hedge itself when implemented correctly, plus the
bond's \term{coupon payments} and the \term{interest earned on the short sale
proceeds}.
\item \textbf{Costs:} \term{borrowing fees} paid to the stock lender and the
\term{dividends} that must be paid to the lender on the borrowed shares.
\item \textbf{Monetising volatility:} because the relationship between the stock
price and the convertible's price is \term{nonlinear}, stock price changes affect
the hedge asymmetrically. Whether the stock rises or falls, the gain or loss on
the short stock position tends to \emph{exceed} the corresponding loss or gain on
the convertible. This is the same effect as convexity in the bond price--yield
relationship.
\item \textbf{Cheapness:} a further source of return arises when the convertible
is itself undervalued, for example when the embedded call is priced below fair
value.
\end{itemize}

\begin{imexbox}[title={Hedged against unhedged, on a $\pm$5\% move}]
A convertible bond is bought at 101.375\% of a \$1,000 par value, so
\term{\$1,013.75}. The stock trades at \$41.54, the convertible's delta is 53\%,
the conversion ratio is 21.2037 shares, and the annual coupon is \$2.50. A
convertible bond pricing model revalues the bond at \$1,038.07 if the stock rises
5\% and \$991.78 if it falls 5\%.

\medskip
\textbf{Unhedged (long the convertible only):}
\begin{align*}
\text{stock } +5\%: &\quad (1{,}038.07 - 1{,}013.75) + 2.50 = +\$26.82 \\
\text{stock } -5\%: &\quad (\phantom{0}991.78 - 1{,}013.75) + 2.50 = -\$19.47
\end{align*}

\textbf{Hedged (short the underlying stock against the bond):} when the stock
declines by 5\%, the arbitrage position produces a \term{profit of \$1.37},
against a \emph{loss} of \$19.47 in the unhedged long-only position. When the
stock rises, the hedged gain is smaller than the unhedged one.
\end{imexbox}

\begin{imtrapbox}[breakable=false,title={What the hedge does to the payoff}]
The hedge does not simply reduce the size of the outcome. It \term{changes the
sign} on the downside: a 5\% fall goes from a \$19.47 loss to a \$1.37 gain. The
price of that is a smaller gain on the upside. Read the objective's verb
carefully --- ``describe the mechanics'' means this asymmetry, not the definition
of a convertible. \emph{Note also that the \$19.47 is not the price change alone:
the price falls \$21.97 and the \$2.50 coupon is added back.} Omitting the coupon
is the standard slip and lands on a real number.
\end{imtrapbox}

\figwrap{%
\begin{tikzpicture}
\begin{axis}[frmaxis,
  width=11.0cm, height=5.4cm,
  ybar, bar width=17pt,
  symbolic x coords={stock $-5\%$, stock $+5\%$},
  xtick=data, x tick label style={font=\scriptsize\sffamily},
  enlarge x limits=0.45,
  ymin=-26, ymax=34,
  ylabel={net gain or loss (\$ per bond)},
  legend style={font=\scriptsize\sffamily, draw=FRMRule, fill=white,
    at={(0.5,-0.24)}, anchor=north, legend columns=-1, column sep=10pt},
]
\addplot[draw=FRMRed,fill=PaleRed,area legend] coordinates
  {(stock $-5\%$,-19.47) (stock $+5\%$,26.82)};
\addlegendentry{unhedged long convertible}
\addplot[draw=FRMGreen,fill=PaleGreen,area legend] coordinates
  {(stock $-5\%$,1.37) (stock $+5\%$,8.0)};
\addlegendentry{hedged convertible arbitrage}
\draw[black!45, thin] (axis cs:stock $-5\%$,0) -- (axis cs:stock $+5\%$,0);
\node[font=\scriptsize\sffamily,color=black!70,anchor=north]
  at (axis cs:stock $-5\%$,-20.2) {$-19.47$};
\node[font=\scriptsize\sffamily,color=FRMGreen,anchor=west,fill=white,inner sep=1.5pt]
  at (axis cs:stock $-5\%$,6.5) {$+1.37$};
\node[font=\scriptsize\sffamily,color=black!70,anchor=south]
  at (axis cs:stock $+5\%$,27.4) {$+26.82$};
\end{axis}
\end{tikzpicture}}%
{Why the trade is worth putting on. The unhedged position is a leveraged bet on
the stock: a 5\% move either way produces a swing of roughly \$20 to \$27 on a
\$1,013.75 bond. The hedge cuts the upside and, more importantly, \emph{turns the
downside positive}. Only the two red bars and the hedged downside figure of
\$1.37 are stated in the source; the hedged upside bar is drawn to scale to show
the shape, not as a reported figure.}

\subsection{Relative value arbitrage}

Relative value arbitrage seeks pricing inefficiencies across asset classes.

\begin{enumerate}
\item \term{Pairs trading.} Two highly correlated firms in the same industry. If
prices diverge without fundamental reason, go long the undervalued and short the
overvalued; profit when the spread reverts to its historical relationship.
\item \term{Statistical arbitrage.} Extends pairs trading across many securities
using factor models --- many small long--short positions identifying temporary
mispricings. One source of profit is \term{liquidity provision}: when large
institutional investors sell big blocks, prices may fall temporarily, and the fund
buys the undervalued stock, shorts a related stock to hedge market risk, and holds
until prices normalise.
\item \term{Index arbitrage.} If index futures are mispriced relative to the
underlying stocks, buy the undervalued side and short the overvalued one. Since
trading every constituent is impractical, algorithms select an optimal subset
considering costs and mispricing. Firms with existing stock holdings have an
advantage, because they can source shares internally at lower cost.
\end{enumerate}

% ------------------------------------------------------------------
\lo{IM-8 c}{Describe the mechanics and risks of event-driven strategies such as merger arbitrage.}

Event-driven strategies focus on specific corporate events that create pricing
discrepancies and market inefficiencies.

\renewcommand{\arraystretch}{1.2}
\noindent\begin{tabularx}{\textwidth}{@{}>{\raggedright\arraybackslash}p{2.7cm}X@{}}
\toprule
\textbf{\sffamily\small Category} & \textbf{\sffamily\small Opportunities} \\
\midrule
Strategic \newline (hard catalysts) & Risk arbitrage; evaluating a firm's strategic alternatives and taking a position; potential spin-offs of divisions or business units; taking positions in likely takeover targets; proxy contests or activist shareholder campaigns; \term{stub trades}, which isolate and profit from the discount between a holding company's share price and the value of its underlying assets. \\
Financial & Anticipating price movements based on credit re-ratings or liquidity events; changes in accounting methods; financial recapitalizations; primary offerings of debt or equity; bankruptcy reorganizations. \\
Operational & Synergies from mergers; corporate turnaround or restructuring opportunities; senior management changes. \\
Legal / regulatory & Positions based on litigation outcomes, on regulatory changes, or on anticipated legislation. \\
Technical & \term{Broken} risk arbitrage situations; secondary offerings of equity or equity-linked notes. \\
\bottomrule
\end{tabularx}
\renewcommand{\arraystretch}{1}
\figcap{Five categories of event-driven opportunity, none of which appears in the
source notes' text layer --- the entire table is an image on the page. Note the
last row: \emph{broken} risk arbitrage is its own opportunity set, which is what
separates a technical catalyst from a strategic one.}

\subsection{Activist strategy}

Taking a stake in a company and actively pushing management to make changes that
increase shareholder value --- governance changes, restructuring, cost cuts,
dividends or buybacks.

\subsection{Merger arbitrage}

The goal of merger arbitrage, also known as \term{risk arbitrage}, is to profit
from the spread between the acquirer's offer price and the target company's stock
price following the announcement of an acquisition or merger.

\begin{imkeybox}[title={The trade, and both sides of the outcome}]
\textbf{Strategy:} buy (long) the target company's stock; short the acquirer's
stock in stock-for-stock deals.

\smallskip
\textbf{If the deal closes --- profit sources:} the arbitrage spread (target price
against final offer price); dividends on the target stock; interest from the short
position, net of costs; and the chance of a higher competing bid.

\smallskip
\textbf{If the deal fails --- risks:} the target stock falls, often back to its
pre-announcement level; and the acquirer's stock may \emph{rise}, producing a loss
on the short leg. \term{Both legs lose at once.}
\end{imkeybox}

\begin{imexbox}[title={Constructing the merger arbitrage trade}]
Arielle Way, Inc.\ (AWI) has offered to acquire Cockrell Industries (CI) for two
shares of AWI stock for each share of CI. Before the announcement AWI traded at
\$52 and CI at \$74, so the offer represents roughly a 40\% premium
$\left(\tfrac{2 \times 52}{74} - 1 = 40.5\%\right)$. Neither firm pays dividends
and the transaction is expected to close within two months. Following the
announcement AWI trades at \$50 and CI at \$95.

\medskip
\textbf{The trade:} buy 100 shares of CI at \$95; sell short 200 shares of AWI at
\$50.

\medskip
\textbf{Anticipated result:} when the transaction closes, the short position
closes, for a profit of
$\dfrac{200(50) - 100(95)}{100} = \mathbf{\$5}$ per CI share. If competing bids
emerge, CI's share price may rise further while AWI's may decline, potentially
increasing profits.
\end{imexbox}

\begin{imexbox}[title={The annualised expected return on a merger arbitrage}]
\[ \text{annualised expected return} = \frac{(C \times G) - \left[(1-C) \times L\right]}{Y \times P} \]
\vk{$C$ is the probability of success; $G$ the expected gain if it succeeds; $L$
the expected loss if it fails; $Y$ the expected time to closing in years; and $P$
the current price of the security.}

\medskip
Global Semiconductors makes a tender offer to shareholders of Local
Semiconductors at \$12 per share. Before the announcement Local traded at \$9;
it rose immediately to \$10.50. The deal is expected to close in six months with
an estimated 90\% probability of completion.

\medskip
$C = 0.90$; $G = 12 - 10.50 = \$1.50$; $L = 10.50 - 9 = \$1.50$; $P = \$10.50$;
$Y = 6/12 = 0.50$.
\[ \text{expected annualised return}
 = \frac{0.90(1.50) - 0.10(1.50)}{0.50 \times 10.50}
 = \frac{1.20}{5.25} = \mathbf{22.86\%} \]
\end{imexbox}

\begin{imtrapbox}[breakable=false,title={Neither of these examples is in the text layer}]
Both worked cases above sit entirely inside images on pages that return fifteen
characters of extractable text. The annualised-return formula in particular
appears nowhere in the written notes. Two details of it are corruptible: $L$ is
measured back to the \term{pre-announcement} price, not to zero; and $P$ is the
\term{current} price of \$10.50, not the offer price of \$12 or the
pre-announcement \$9.
\end{imtrapbox}

% ------------------------------------------------------------------
\lo{IM-8 d}{Explain the mechanics of a distressed securities strategy, and describe challenges, risks and considerations that may arise when implementing this strategy.}

A distressed securities strategy involves taking speculative positions in the
securities of firms that are \term{either entering, in, or emerging from
bankruptcy}. Examples of potential returns include:

\begin{enumerate}[label=\alph*)]
\item Buying \term{equity} and benefiting from price recovery after
restructuring, or purchasing non-investment-grade bonds at discounted prices, as
many institutions are \emph{restricted} from holding such debt.
\item Banks selling \term{distressed loans} to improve their balance sheets,
allowing investors to acquire them at attractive valuations.
\item \term{Trade claims} purchased at discounts from holders unwilling to stay
through the restructuring process.
\item Exploiting \term{mispricing within the capital structure} --- for example
long senior securities and short junior ones, expecting priority-based recovery
differences.
\end{enumerate}

\begin{imkeybox}[title={Two investor approaches}]
\term{Active:} participate in creditor committees, influence restructuring
decisions, oversee the workout process, or engage in legal action to improve
recovery. \quad
\term{Passive:} avoid involvement and focus on less complex distressed
opportunities with clearer outcomes.
\end{imkeybox}

\begin{imexbox}[title={A capital-structure distressed trade}]
Naples Ocean Views (NOV), a real estate holding company, may be forced to file for
bankruptcy. The firm has \$1 billion of 5\% debt maturing in 10 years, currently
trading at a steep discount of \$30 per \$100 of face value. Bankruptcy is
expected within two years. NOV's stock trades at \$3 per share. In liquidation
the bonds are estimated to be worth \$35 per \$100 of face value. An investor
buys one NOV bond of \$1,000 face at \$300 and shorts 100 shares at \$3. The firm
is liquidated in two years.

\medskip
\textbf{Gain on the bond position:} coupons of 5\% on \$1,000 for two years
($\$50$ per year, \$100 total), plus \$350 in liquidation value, for total
inflows of \$450.
\[ \$450 - \$300 = \mathbf{\$150} \]

\textbf{Gain on the equity short:} short sale proceeds are \$300 and the stock
becomes worthless, so the shares never have to be bought back.
\[ \mathbf{\$300} \]

\textbf{Total gain} $= \$150 + \$300 = \mathbf{\$450}$.
\end{imexbox}

\begin{imtrapbox}[breakable=false,title={Why this is the capital-structure trade in miniature}]
The bond is bought at \$30 and recovers \$35 --- a modest \$50 of principal gain,
with the coupons doing as much work again. The equity is the leg that pays: it
goes to zero, and the entire \$300 of short proceeds is kept. \term{Priority is
the whole trade.} Note the two easy slips: the coupons run for \emph{two} years,
not one, and the \$35 recovery is per \$100 of \emph{face}, so \$350 on a \$1,000
bond, not \$35.
\end{imtrapbox}

\subsection{Key risks and challenges}

High uncertainty in valuation, due to unstable cash flows and complex capital
structures, compounded by illiquidity and by the length and unpredictability of
the legal process.

\begin{imtrapbox}[title={Consolidated trap summary --- IM-8}]
\begin{itemize}
\item \term{Source, all objectives.} No GARP book covers this reading. Everything
gap-filled here is one rung below a GARP source.
\item \term{Polarity, IM-8 a.} Dedicated short faces a \emph{natural headwind}
because markets rise over time, which is why the strategy needs unusual expertise
rather than merely the opposite skill to going long.
\item \term{Sibling, IM-8 a.} Global macro uses the \emph{most} liquid
instruments; emerging markets accepts illiquidity. Both are macro.
\item \term{Definition, IM-8 b.} Convertible arbitrage is long the bond, short the
stock. The hedge must be \emph{delta} hedged and continuously rebalanced, because
delta moves with the share price.
\item \term{Sign, IM-8 b.} The hedge turns a 5\% downside from a \$19.47 loss into
a \$1.37 gain. It does not merely shrink the loss.
\item \term{Intermediate result, IM-8 b.} The \$19.47 includes the \$2.50 coupon
added back to a \$21.97 price decline. Dropping the coupon lands on \$21.97.
\item \term{Role, IM-8 c.} In a stock-for-stock deal you are long the
\emph{target} and short the \emph{acquirer}. If the deal breaks, both legs lose.
\item \term{Formula, IM-8 c.} In the annualised merger return, $L$ runs back to
the \emph{pre-announcement} price and $P$ is the \emph{current} price.
\item \term{Intermediate result, IM-8 d.} A \$35 recovery per \$100 of face is
\$350 on a \$1,000 bond. Coupons run for the full two years.
\end{itemize}
\end{imtrapbox}

% ==================================================================
